\reading{IM-4}{Portfolio Construction}{DEFER}{1}

\srcnote{Richard Grinold and Ronald Kahn, \textit{Active Portfolio Management: A
Quantitative Approach for Producing Superior Returns and Controlling Risk}, 2nd
Edition (New York, NY: McGraw-Hill, 2000), Chapter 14. Printed as Chapter 4 in
the GARP source volume.}

\begin{imkeybox}[title={The nine objectives in this reading}]
\begin{enumerate}[label=\textbf{\alph*.}]
\item Describe the inputs to the portfolio construction process and explain challenges faced when using these inputs.
\item Evaluate the motivation for and the methods used for refining alphas in the implementation process.
\item Describe neutralization and the different approaches used for refining alphas to be neutral.
\item Explain the implications of transaction costs on portfolio construction.
\item Describe practical issues in portfolio construction, including the determination of an appropriate risk aversion, aversions to specific risks, and proper alpha coverage.
\item Describe portfolio revisions and rebalancing, and analyze the tradeoffs between alpha, risk, transaction costs, and time horizon.
\item Determine the optimal no-trade region for rebalancing with transaction costs.
\item Evaluate the strengths and weaknesses of the following portfolio construction techniques: screens, stratification, linear programming, and quadratic programming.
\item Describe dispersion, explain its causes, and describe methods for controlling forms of dispersion.
\end{enumerate}
\end{imkeybox}

\begin{imtrapbox}[breakable=false,title={Nine objectives at SRC 1}]
This is the second-largest objective count in the area sitting on the lowest
historical yield in it. The volume covers it in full; the tag tells you where to
spend revision time, not what to omit.
\end{imtrapbox}

% ------------------------------------------------------------------
\lo{IM-4 a}{Describe the inputs to the portfolio construction process and explain challenges faced when using these inputs.}

\figwrap{%
\begin{tikzpicture}[
  inp/.style={draw=FRMNavy,fill=PaleNavy,rounded corners=1pt,inner sep=3.5pt,
              font=\scriptsize\sffamily,text=FRMNavy,align=center},
  hub/.style={draw=FRMRed,fill=PaleRed,rounded corners=2pt,inner sep=5pt,
              font=\small\sffamily\bfseries,text=FRMRed,align=center},
  tec/.style={draw=FRMRule,fill=white,rounded corners=1pt,inner sep=3.5pt,
              font=\scriptsize\sffamily,align=center,text width=30mm},
  ar/.style={-{Stealth[length=4pt]},FRMGold,line width=0.9pt}]
\node[inp] (i1) at (-5.4,2.1) {Current portfolio};
\node[inp] (i2) at (-2.7,2.7) {Alphas};
\node[inp] (i3) at (0.2,2.7) {Covariance};
\node[inp] (i4) at (3.0,2.1) {Active risk aversion};
\node[inp] (i5) at (-4.2,1.15) {Transactions cost};
\node[hub] (h) at (-1.2,0.2) {Portfolio construction};
\foreach \n in {i1,i2,i3,i4,i5} \draw[ar] (\n) -- (h);
\node[tec] (t1) at (-4.4,-1.5) {Screens};
\node[tec] (t2) at (-1.2,-1.5) {Stratification};
\node[tec] (t3) at (2.0,-1.5) {Linear programming};
\node[tec] (t4) at (5.2,-1.5) {Quadratic programming};
\draw[ar] (h) -- node[right,font=\scriptsize\sffamily,text=black!60,
  fill=white,inner sep=1.5pt]{techniques} (-1.2,-1.0);
\draw[FRMGold,line width=0.9pt] (-4.4,-1.0) -- (5.2,-1.0);
\foreach \n in {t1,t2,t3,t4} \draw[ar] (\n.north|-{(0,-1.0)}) -- (\n.north);
\end{tikzpicture}}%
{The whole reading on one page. Five inputs feed the construction problem; four
techniques solve it. Every later objective in this reading is either a way of
cleaning one of the five inputs (refining alphas, neutralisation, choosing a risk
aversion) or a consequence of one of them being costly (transaction costs, the
no-trade region, dispersion). Redrawn from the source notes' own sketch.}

We use mean-variance analysis to find the best mix based on these inputs.
However, real-world limitations can restrict portfolio weights.

\gapbadge
\srcnote{The spine's objective has two halves --- the inputs \emph{and} the
challenges faced when using them --- and the notes' heading carries only the
first. The challenges are written below from GARP Chapter 4.}

\begin{imgapbox}[title={What the objective asks for}]
Not a list of five nouns, but which of the five you can actually trust.
\end{imgapbox}

\begin{imkeybox}[title={Only one input is reliable}]
\begin{itemize}
\item \term{The current portfolio} is the only input we can measure with near
certainty.
\item \term{Alphas} are often unreasonable and subject to hidden biases.
\item \term{Covariances} and \term{transactions cost} estimates are noisy; we
hope they are unbiased, but we know they are not measured with certainty.
\item \term{Active risk aversion} is not certain either. Most active managers
have a target level of active risk that must be made consistent with an active
risk aversion.
\end{itemize}
Much of what follows in this reading is an indirect method of coping with noisy
inputs.
\end{imkeybox}

% ------------------------------------------------------------------
\lo{IM-4 b}{Evaluate the motivation for and the methods used for refining alphas in the implementation process.}

With alpha analysis, the alphas can be adjusted so that they are in line with the
manager's desires for risk control and anticipated sources of value added. The
three methods are:

\begin{enumerate}
\item \term{Scale the alphas.}
\item \term{Trim alpha outliers.}
\item \term{Neutralization.}
\end{enumerate}

\begin{imkeybox}
Using refined alphas and then performing optimization can achieve the same goal
as a constrained optimization approach, but has the advantage of focusing on the
alpha inputs and the effects of individual constraints on portfolio returns.
\end{imkeybox}

% ------------------------------------------------------------------
\lo{IM-4 c}{Describe neutralization and the different approaches used for refining alphas to be neutral.}

\begin{imdefbox}[title={Neutralization}]
Neutralization is the removal of biases or undesirable bets from our alphas.
\end{imdefbox}

\begin{itemize}
\item \term{Benchmark-neutral alphas.} Means that the benchmark has zero alpha.
If our initial alphas imply an alpha for the benchmark, the neutralization
process recenters the alphas to remove the benchmark alpha.
\item \term{Cash-neutral alphas.} Alphas will not lead to any active cash
position.
\item \term{Risk-factor-neutral alphas.} If the manager does not have any ability
to forecast the risk factors, we can use risk-factor-neutral alphas. Once
neutralized, the alphas of the risk factors will be zero.
\end{itemize}

% ------------------------------------------------------------------
\lo{IM-4 d}{Explain the implications of transaction costs on portfolio construction.}

\begin{itemize}
\item Trading fees and spreads eat into your returns compared to the benchmark.
\item Unlike alphas, which are uncertain but predictable, \term{transaction costs
are more certain but have some variability}.
\item They convert a one-dimensional problem --- the trade-off between alpha and
active risk --- into a \term{two-dimensional problem}: alpha, active risk, and
transaction costs.
\item Constructing, rebalancing and liquidating all incur transaction costs.
\end{itemize}

\begin{imtrapbox}[breakable=false,title={Costs are a point in time, returns are a period}]
We must be careful to compare transactions costs incurred at a \emph{point in
time} with returns and risk realized over a \emph{period of time}. That is why
the costs have to be amortised over the horizon across which the alpha is
expected to be realised --- and why the horizon is inseparable from the
alpha/risk/cost trade-off in IM-4 f.
\end{imtrapbox}

% ------------------------------------------------------------------
\lo{IM-4 e}{Describe practical issues in portfolio construction, including the determination of an appropriate risk aversion, aversions to specific risks, and proper alpha coverage.}

\subsection{Determination of risk aversion}

If we have better intuition about our information ratio and our desired amount of
active risk, we can back out the risk aversion:

\begin{imfmlbox}
\[ \lambda_A = \frac{\text{IR}}{2\,\sigma_A} \]
\vk{$\lambda_A$ is the active risk aversion, IR the information ratio, and
$\sigma_A$ the desired active risk. \term{Percentages, not decimals}: verify what
the optimizer expects.}
\end{imfmlbox}

\begin{imexbox}[title={Backing out a risk aversion}]
If our information ratio is $0.5$ and we desire $5\%$ active risk:
\[ \lambda_A = \frac{0.5}{2 \times 5} = 0.05 \]
\end{imexbox}

\subsection{Incorporation of specific risk aversion}

Beyond overall risk, some investors dislike specific risks, for example a sector
mismatch with the benchmark. A high aversion to specific risk reduces bets on any
one stock.

\subsection{Proper alpha coverage}

What about stocks not in the benchmark, or without alpha forecasts?

\begin{itemize}
\item \term{Stocks outside the benchmark:} assign a zero benchmark weight but
allow for an active weight, to capture potential alpha.
\item \term{Stocks without alpha forecasts:} infer adjusted alphas based on those
with forecasts.
\end{itemize}

% ------------------------------------------------------------------
\lo{IM-4 f}{Describe portfolio revisions and rebalancing, and analyze the tradeoffs between alpha, risk, transaction costs, and time horizon.}

\begin{itemize}
\item If a manager knows how to make the correct trade-off between expected
active return, active risk and transactions costs, \term{frequent revision will
not present a problem}.
\item If the manager is unsure of her ability to correctly specify alphas, active
risk and transactions costs, she may resort to \term{less frequent revision as a
safeguard}.
\end{itemize}

% ------------------------------------------------------------------
\lo{IM-4 g}{Determine the optimal no-trade region for rebalancing with transaction costs.}

\begin{imfmlbox}
\[ \text{MCVA}_n = \alpha_n - 2 \times \lambda_A \times \psi \times \text{MCAR}_n \]
\vk{$\text{MCVA}_n$ is stock $n$'s marginal contribution to value added --- how
risk-adjusted alpha changes as the holding of the stock is increased, with an
offsetting decrease in cash. $\text{MCAR}_n$ is its marginal contribution to
active risk, the rate at which active risk changes as we add more of stock $n$.
$\lambda_A$ is the active risk aversion and $\psi$ the active risk.}
\end{imfmlbox}

We can decide whether to trade by comparing the MCVA to the transactions costs.
\term{As long as $\text{MCVA}_n$ stays within the negative cost of selling and
the cost of purchase, the portfolio will remain optimal}, and we should not react
to new information.

\begin{imfmlbox}
\[ -\text{SC}_n \;\le\; \text{MCVA}_n \;\le\; \text{PC}_n \qquad\Longrightarrow\qquad \text{do not trade} \]
\vk{$\text{PC}_n$ is the purchase cost and $\text{SC}_n$ the sales cost for stock
$n$. The band is \emph{asymmetric} whenever the two costs differ.}
\end{imfmlbox}

\gapbadge
\srcnote{The notes state the band rule without ever putting numbers through it.
The worked case below is GARP's own (Chapter 4), and it is what makes the
asymmetry visible.}

\begin{imexbox}[title={Working the band}]
Take $\text{PC}_n = 0.50\%$ and $\text{SC}_n = 0.75\%$, so the no-trade band runs
from $-0.75\%$ to $+0.50\%$.
\begin{itemize}
\item If $\text{MCVA}_n = 0.80\%$, it exceeds the purchase cost. Buying yields a
net benefit of $0.80\% - 0.50\% = \mathbf{0.30\%}$.
\item If $\text{MCVA}_n = -1.30\%$, it is below the negative of the sales cost.
Decreasing the holding saves $1.30\%$ at the margin against a $0.75\%$ cost, for
a net benefit of $1.30\% - 0.75\% = \mathbf{0.55\%}$.
\end{itemize}
\end{imexbox}

\figwrap{%
\begin{tikzpicture}
\begin{axis}[frmaxis,
  width=11.4cm, height=5.8cm,
  xmin=0, xmax=10, ymin=-1.75, ymax=1.25,
  xtick=\empty,
  ylabel={$\text{MCVA}_n$ (\%)},
  ytick={-1.30,-0.75,0,0.50,0.80},
  yticklabels={$-1.30$,$-0.75$,$0$,$0.50$,$0.80$},
]
\fill[PaleGreen] (axis cs:0,-0.75) rectangle (axis cs:10,0.50);
\draw[FRMGreen, thick] (axis cs:0,0.50) -- (axis cs:10,0.50);
\draw[FRMGreen, thick] (axis cs:0,-0.75) -- (axis cs:10,-0.75);
\draw[black!45, thin] (axis cs:0,0) -- (axis cs:10,0);
\addplot[only marks, mark=*, mark size=2.2pt, FRMRed] coordinates {(3,0.80) (7,-1.30)};
\node[font=\scriptsize\sffamily, fill=white, inner sep=1.5pt, anchor=west,
      text=FRMRed] at (axis cs:3.3,0.80) {buy: net $+0.30\%$};
\node[font=\scriptsize\sffamily, fill=white, inner sep=1.5pt, anchor=west,
      text=FRMRed, anchor=east] at (axis cs:6.7,-1.30) {sell: net $+0.55\%$};
\node[font=\scriptsize\sffamily, fill=white, inner sep=1.5pt, anchor=west,
      text=FRMGreen] at (axis cs:0.3,-0.14) {no-trade region};
\node[font=\scriptsize\sffamily, anchor=west, text=black!55]
  at (axis cs:8.6,0.60) {$+\text{PC}_n$};
\node[font=\scriptsize\sffamily, anchor=west, text=black!55]
  at (axis cs:8.6,-0.65) {$-\text{SC}_n$};
\end{axis}
\end{tikzpicture}}%
{The no-trade band, drawn with GARP's own costs of $0.50\%$ to purchase and
$0.75\%$ to sell. Note that the band is not symmetric about zero: because selling
costs more than buying, the portfolio tolerates a larger \emph{negative} MCVA
before it is worth acting. New information that moves MCVA anywhere inside the
shaded region should be ignored --- which is the objective's real content, since
the intuitive answer is that any new information warrants a trade.}

% ------------------------------------------------------------------
\lo{IM-4 h}{Evaluate the strengths and weaknesses of the following portfolio construction techniques: screens, stratification, linear programming, and quadratic programming.}

\renewcommand{\arraystretch}{1.25}
\noindent\begin{tabularx}{\textwidth}{@{}>{\raggedright\arraybackslash}p{2.6cm}XX@{}}
\toprule
\textbf{\sffamily\small Technique} & \textbf{\sffamily\small Strengths} & \textbf{\sffamily\small Weaknesses} \\
\midrule
Screens & Simple to use; allows filtering based on alpha or other criteria. \emph{Example:} selecting top alpha stocks for your portfolio from 100 stocks. & Can be overly simplistic; ignores diversification. \\
Stratification & Ensures the portfolio reflects the benchmark's risk profile. \emph{Example:} selecting stocks by category (size, sector). & Limits flexibility; might miss alpha opportunities outside categories. \\
Linear programming & More sophisticated than screens; considers multiple risk factors (size, beta). Can incorporate transaction costs and position size limits. & May result in a portfolio very close to the benchmark, potentially missing alpha. \\
Quadratic programming & Theoretically the best method; uses alphas, risks and transactions costs for an optimal portfolio. Allows various constraints to be incorporated. & Highly sensitive to estimation errors, especially in covariances. Requires accurate data for all assets: 400 stocks lead to 79,800 covariance estimates. \\
\bottomrule
\end{tabularx}
\renewcommand{\arraystretch}{1}
\figcap{The four techniques, in increasing order of sophistication and of
vulnerability. The progression is not a ranking of which to use: the last row's
strength and weakness are the same property. Quadratic programming uses every
input, which is exactly why it is the most exposed to the noise in those inputs
that IM-4 a warned about.}

\begin{imkeybox}[title={Where 79,800 comes from}]
A covariance matrix over $N$ assets has $N(N-1)/2$ distinct pairwise covariances.
For $N = 400$:
\[ \frac{400 \times 399}{2} = 79{,}800 \]
plus 400 variances on the diagonal, for 80,200 distinct entries in total. Every
one of them is estimated, and every one carries error.
\end{imkeybox}

% ------------------------------------------------------------------
\lo{IM-4 i}{Describe dispersion, explain its causes, and describe methods for controlling forms of dispersion.}

\begin{imdefbox}[title={Dispersion}]
The difference between the maximum return and the minimum return for separate
account portfolios. It is a measure of how an individual client's portfolio may
differ from the manager's reported composite returns.
\end{imdefbox}

\subsection{Sources of dispersion}

\begin{itemize}
\item \term{Client-driven:} clients impose different constraints.
\item Lack of attention to separate accounts.
\end{itemize}

\subsection{Controlling dispersion}

If transactions costs were zero, rebalancing all the separate accounts so that
they hold exactly the same assets in the same proportions would have no cost.
Dispersion would disappear, at no cost to investors.

With transactions costs, \term{managers should reduce dispersion only until
further reduction would substantially lower returns on average}, because much
higher transactions costs would be incurred.

\begin{imtrapbox}[title={Consolidated trap summary --- IM-4}]
\begin{itemize}
\item \term{Role, IM-4 a.} Only the \emph{current portfolio} is measured with
near certainty. Alphas, covariances, transaction costs and risk aversion are all
uncertain. A question that calls any of the other four reliable is corrupting
this.
\item \term{Formula, IM-4 e.} $\lambda_A = \text{IR}/(2\sigma_A)$, and the inputs
are in \emph{percent}. IR $0.5$ with $5\%$ active risk gives $0.05$, not $0.005$.
\item \term{Polarity, IM-4 g.} Trade when MCVA lies \emph{outside} the band. New
information inside the band is ignored.
\item \term{Sign, IM-4 g.} The band is $[-\text{SC}_n, +\text{PC}_n]$: the sales
cost enters \emph{negatively}. Swapping the two ends, or making the band
symmetric, is the standard corruption.
\item \term{Sibling, IM-4 c.} Benchmark-neutral sets the \emph{benchmark's} alpha
to zero; risk-factor-neutral sets the \emph{factors'} alphas to zero;
cash-neutral removes the active \emph{cash} position. Three different things
being set to zero.
\item \term{Definition, IM-4 d.} Alphas are uncertain but predictable;
transaction costs are more certain but variable. The asymmetry runs in the
direction people do not expect.
\item \term{Intermediate result, IM-4 h.} $79{,}800$ is the count of distinct
\emph{pairwise covariances}, $N(N-1)/2$. The full matrix has $80{,}200$ distinct
entries once the variances are included.
\item \term{Scope, IM-4 i.} Dispersion is about \emph{separate accounts} against
the composite, not about tracking error against a benchmark.
\end{itemize}
\end{imtrapbox}

% ==================================================================
